\subsubsection{Summarization}
\label{sec:human-evaluation-summarization}
For the two summarization scenarios (\cnndm, \xsum), we conduct human evaluations for summary quality, with a special interest in the faithfulness of summaries.
To explore the results further, including the model generations, see  \benchmarkUrlPrefix{summarization}.

\paragraph{Evaluation details.}
To evaluate models in general, we conducted \numruns runs with \numevalinstances evaluation instances (for both \cnndm and \xsum).
For the purposes of human evaluation, we used 100 instances sampled from the full \numevalinstances instances for the first run, and used the same 100 instances for all models (for both \cnndm and \xsum).
In addition to human evaluation for the few-shot performance of the aforementioned six language models, as a one-off analysis in light of the findings of \citet{goyal2022news} on the zero-shot performance of language models, we evaluated \gptdavinci, \gptcurie, \instructdavinci, and \instructdavinci zero-shot.

In addition to human evaluation for the six language models under few-shot conditions and the four language models under zero-shot conditions, we evaluated two state-of-the-art models that were extensively fine-tuned on the associated datasets: Pegasus~\citep{zhang2020pegasus} and BRIO~\citep{liu-etal-2022-brio}.
And given  known issues in \cnndm and \xsum \citep[][]{bommasani2020intrinsic, gehrmann2022repairing, reiter2022summarization}, we conduct human evaluation on the \textit{reference} summaries.
In total, this means we conduct human evaluation on 13 sets of summaries for each dataset: six sets generated by language models adapted through few-shot prompting, four sets generated by language models adapted through zero-shot prompting (just instructions), two sets generated by state-of-the-art models fine-tuned for the respective datasets, and the official human-authored reference summaries. 

\paragraph{Annotation guidelines.}
We ask the annotators to evaluate each summary for three aspects: \textit{faithfulness}, \textit{coherence}, and \textit{relevance}, following the guidelines of \citet{fabbri-etal-2021-summeval}.\footnote{We omit evaluating fluency because we find language model outputs are mostly fluent.}
We define a summary to be faithful if ``all the information expressed by the summary can be inferred from the article'' and solicit binary decisions from the annotators.
We define a summary to be relevant if the summary ``includes only important information from the source document'' and coherent if the summary ``organizes the relevant information into a well-structured summary''.
For relevance and coherence, we ask the annotators to annotate on a 1--5 Likert scale.

We recruit annotators from Amazon Mechanical Turk, compensating them at California minimum wage of \$15.00/hr using conservative estimates of annotation time, consistent with best-practices for crowd-sourcing labor \citep{whiting2019fair}.\footnote{We estimated the time it would take workers by performing the task ourselves (it took on average 1 minute per HIT). We added a buffer of $25\%$, which would result in $31$ cents per HIT at the minimum wage (we rounded it up to $35$ cents).}
Each model summary was evaluated by three annotators; we report results based on their average score for each summary for brevity.\footnote{All underlying annotations will be made available to facilitate future analysis.}

\begin{table*}[htp]
\resizebox{\textwidth}{!}{
\begin{tabular}{llrrrrrr}
\toprule
 &  & \multicolumn{3}{c}{\cnndm} & \multicolumn{3}{c}{\xsum} \\
Setting & Models & Faithfulness & Coherence & Relevance & Faithfulness & Coherence & Relevance \\
\midrule
\multirow[c]{4}{*}{Zero-shot language models} & \gptcurie & 0.29 & 1.77 & 1.93 & 0.77 & 3.16 & 3.39 \\
 & \gptdavinci & 0.76 & 2.65 & 3.50 & 0.80 & 2.78 & 3.52 \\
 & \instructcurie & 0.97 & 4.24 & 4.59 & 0.96 & 4.27 & 4.34 \\
 & \instructdavinci & 0.99 & 4.15 & 4.60 & 0.97 & 4.41 & 4.28 \\
\midrule
\multirow[c]{6}{*}{Five-shot language models} & \anthropic & 0.94 & 3.88 & 4.33 & 0.70 & 4.77 & 4.14 \\
 & \coherexl & 0.99 & 3.42 & 4.48 & 0.63 & 4.79 & 4.00 \\
 & \glm & 0.94 & 3.69 & 4.24 & 0.74 & 4.72 & 4.12 \\
 & \optonesevenfive & 0.96 & 3.64 & 4.33 & 0.67 & 4.80 & 4.01 \\
 & \gptdavinci & 0.99 & 3.95 & 4.34 & 0.69 & 4.69 & 4.03 \\
 & \instructdavinci & 0.98 & 4.13 & 4.49 & 0.77 & 4.83 & 4.33 \\
\midrule
\multirow[c]{2}{*}{Fine-tuned language models} & Brio & 0.94 & 3.94 & 4.40 & 0.58 & 4.68 & 3.89 \\
 & Pegasus & 0.97 & 3.93 & 4.38 & 0.57 & 4.73 & 3.85 \\
\midrule
Human generated & Reference summaries & 0.84 & 3.20 & 3.94 & 0.37 & 4.13 & 3.00 \\
\bottomrule
\end{tabular}
}
\caption{\textbf{Human evaluation for summarization scenarios.} 
We conduct human evaluation for 13 sets of summaries for both \cnndm and \xsum.}
\label{tab:summarization-human-evaluation}
\end{table*}

\paragraph{Results and discussion.}
\reftab{summarization-human-evaluation} displays the results of our human evaluation for summarization.
First and foremost, we highlight the surprisingly striking finding that the reference summaries are \textbf{low quality}.
In terms of faithfulness, the reference summaries are worse than all other models on \xsum and only outperform the zero-shot non-(instruction-tuned) models on \cnndm. 
In fact, overall we observe that models with greater supervision are less faithful: zero-shot models > few-shot models > finetuned models.

Second, we highlight the importance of instruction-tuning for strong summarization performance.
Consistent with the concurrent work of \citet{goyal2022news}, we observe that zero-shot instruction-tuned models achieve the best summarization accuracy.
However, by evaluating models that both use and don't use instruction tuning, we clarify that instruction tuning is crucial: \gptdavinci performs much worse in terms of all three aspects when compared to \instructdavinci.
Looking into the model behavior, we find that the vanilla GPT-3 models often fail to follow the task instruction, frequently leading to  generations that are either exact repetitions of the source article or irrelevant to the source article.

Third, when we compare the results of human evaluations to automated evaluations, we find the two are \textit{anti-correlated}.
\texttt{ROUGE-2} scores favor fine-tuned models, whereas human judgments consistently prefer few-shot or zero-shot language models. Similarly, we find that automated faithfulness metrics are not reliable for measuring faithfulness of few-shot and zero-shot models. 
For instance, on \xsum, few-shot and zero-shot models score similarly according to automated faithfulness metrics, whereas humans overwhelmingly rate zero-shot models as more faithful.


Put together, our human evaluation presents a worrying state of affairs for summarization in the age of modern language models, especially when using \cnndm and \xsum. 
While these two datasets have served as mainstays in summarization evaluation, often functioning as the standard representatives of summarization in most contexts, we believe they will hold back and, worse, mislead progress in summarization.
We actively encourage the development of high quality summarization evaluation data, and the adoption of such data as mainstream, along with future studies on the quality of summarization metrics as well as the quality of in-context learning and other few-shot methodologies in comparison to fine-tuning.